\uuid{6NFB}
\titre{Module/argument}
\auteur{bodin ; hanani ; mzari}
\organisation{exo7}
\licence{CC BY-NC-SA 4.0}
\datecreate{2020-01-20}
\niveau{L1}
\module{Algèbre}
\chapitre{Nombres complexes}
\sousChapitre{Forme cartésienne, forme polaire}
\difficulte{3}
\theme{Écritures algébrique et géométrique, module, argument}
\exo7qcmid{lille-1.nombres-complexes.015}

\contenu{
\begin{question}
\qclassification{104.01}
\qtags{motcle=module/argument}
Soit $x\in \Rr\backslash \{2k\pi, k \in \Zz\}$, $n \in \Nn^*$,
$S_1= \sum_{k=0}^{n} \cos(kx)$ et $S_2= \sum_{k=0}^{n} \sin(kx)$. Quelles sont les assertions vraies ?
\begin{answers}
    \good{$S_1= \cos (\frac{nx}{2})\cdot  \frac{\sin (\frac{n+1}{2})x}{\sin (\frac{x}{2})}$}

    \bad{$S_1= \sin (\frac{nx}{2}) \cdot  \frac{\sin (\frac{n+1}{2})x}{\sin (\frac{x}{2})}$}

    \good{$S_2=\sin (\frac{nx}{2}) \cdot  \frac{\sin (\frac{n+1}{2})x}{\sin (\frac{x}{2})}$}

    \bad{$S_2= \cos (\frac{nx}{2}) \cdot  \frac{\sin (\frac{n+1}{2})x}{\sin (\frac{x}{2})}$}

\end{answers}
\begin{explanations}
On calcule la somme géométrique $\sum_{k=0}^{n} e^{ikx}= \sum_{k=0}^{n} (e^{ix})^k = \frac{1-e^{i(n+1)x}}{1-e^{ix}}=\frac{e^{i\frac{(n+1)x}{2}}(e^{-i\frac{(n+1)x}{2}}-e^{i\frac{(n+1)x}{2}})}{e^{i\frac{x}{2}}(e^{-i\frac{x}{2}}-e^{i\frac{x}{2}})}= e^{i\frac{nx}{2}}\cdot  \frac{\sin (\frac{n+1}{2})x}{\sin (\frac{x}{2})}$; puis, la partie réelle
et imaginaire de cette somme.
\end{explanations}
\end{question}
}
